\documentclass[12pt]{article}
\usepackage{amsmath}
\usepackage{amssymb}
\usepackage{booktabs}
\usepackage{geometry}
\usepackage{parskip}
\usepackage{hyperref}

\geometry{margin=1in}

\title{Diffractometer Geometry: Problems 1 and 2}
\date{\today}

\begin{document}

\maketitle

\section*{Equipment Overview}

A piece of mechanical equipment consists of two independent stacks of rotary stages.
All rotational axes coincide ideally at a single point of intersection, at which the
object to be examined is mounted.

\subsection*{Reference Frame}

We define a right-handed coordinate system with basis vectors assigned as follows:

\begin{align*}
\hat{x} &= \text{longitudinal (toward equipment, our line of sight)} \\
\hat{y} &= \text{lateral (to our left)} \\
\hat{z} &= \text{vertical (normal to the floor, pointing upward)}
\end{align*}

This satisfies the right-hand rule: $\hat{x} \times \hat{y} = \hat{z}$.

\subsection*{Equipment Description}

The equipment consists of two independent stacks of rotary stages. In the default
orientation (all angles at $0^\circ$), the stacks are described as follows.

\textbf{Stack 1 --- Detector Stack}

\begin{itemize}
  \item \textbf{S1-1}: rotation axis in the vertical direction ($+\hat{z}$).
    Sign of rotation consistent with the coordinate system (right-hand rule about $+\hat{z}$).
  \item \textbf{S1-2}: sits on S1-1. Rotation axis in the lateral direction ($+\hat{y}$).
    Positive rotation takes $+\hat{x}$ toward $+\hat{z}$.
    The detector is mounted on a radial arm on S1-2, directed toward the point of intersection.
\end{itemize}

\textbf{Stack 2 --- Sample Stack}

\begin{itemize}
  \item \textbf{S2-1}: rotation axis in the vertical direction ($+\hat{z}$).
    Axis is colinear with S1-1; the two stages are mechanically independent.
    Sign of rotation consistent with the coordinate system.
  \item \textbf{S2-2}: sits on S2-1. Rotation axis in the lateral direction ($+\hat{y}$).
    Same sign of rotation as S1-2.
  \item \textbf{S2-3}: sits on S2-2. Rotation axis in the longitudinal direction ($+\hat{x}$).
    Sign of rotation consistent with the coordinate system (right-hand rule about $+\hat{x}$).
  \item \textbf{S2-4}: sits on S2-3. Rotation axis in the vertical direction ($+\hat{z}$).
    Sign of rotation consistent with the coordinate system.
\end{itemize}

%-----------------------------------------------------------------------
\section*{Problem 1}

\subsection*{Basis Vector Assignment}

We assign basis vectors to the reference frame as follows:

\begin{align*}
\hat{x} &= \text{longitudinal} \\
\hat{y} &= \text{lateral} \\
\hat{z} &= \text{vertical}
\end{align*}

Right-handedness is confirmed by $\hat{x} \times \hat{y} = \hat{z}$.

\subsection*{Stage Orientations and Sign of Rotation}

The sign of rotation follows the right-hand rule: the thumb points along the
positive axis vector and the fingers curl in the direction of positive rotation.

\textbf{Stack 1 --- Detector Stack}

\begin{center}
\begin{tabular}{llll}
\toprule
Stage & Physical Axis & Axis Vector & Positive Rotation \\
\midrule
S1-1 & vertical & $+\hat{z}$ & CCW viewed from $+\hat{z}$ \\
S1-2 & lateral  & $+\hat{y}$ & $+\hat{x}$ toward $+\hat{z}$ \\
\bottomrule
\end{tabular}
\end{center}

The detector is mounted on a radial arm on S1-2, directed inward toward the sample point.

\textbf{Stack 2 --- Sample Stack}

\begin{center}
\begin{tabular}{llll}
\toprule
Stage & Physical Axis & Axis Vector & Positive Rotation \\
\midrule
S2-1 & vertical     & $+\hat{z}$ & CCW viewed from $+\hat{z}$ \\
S2-2 & lateral      & $+\hat{y}$ & $+\hat{x}$ toward $+\hat{z}$ \\
S2-3 & longitudinal & $+\hat{x}$ & $+\hat{y}$ toward $+\hat{z}$ \\
S2-4 & vertical     & $+\hat{z}$ & CCW viewed from $+\hat{z}$ \\
\bottomrule
\end{tabular}
\end{center}

S2-1 and S1-1 share the same axis of rotation but are mechanically independent.

\subsection*{Computing the Orientation Matrix \textbf{U}}

\textbf{Step 1 --- Define the lab frame.}
The lab frame is fixed with basis $\{\hat{x},\, \hat{y},\, \hat{z}\}$ as assigned above.

\textbf{Step 2 --- Rotation matrix for each stage.}
Each stage $i$ with unit rotation axis $\hat{n}_i$ and rotation angle $\theta_i$
contributes a rotation matrix $\mathbf{R}_i(\theta_i)$ via the Rodrigues formula:

\begin{equation}
  \mathbf{R}_i(\theta_i)
  = \mathbf{I}\cos\theta_i
  + (1 - \cos\theta_i)\,(\hat{n}_i \otimes \hat{n}_i)
  + \sin\theta_i\,[\hat{n}_i]_\times
\end{equation}

where $[\hat{n}]_\times$ is the skew-symmetric cross-product matrix:

\begin{equation}
  [\hat{n}]_\times =
  \begin{pmatrix}
     0   & -n_z &  n_y \\
     n_z &  0   & -n_x \\
    -n_y &  n_x &  0
  \end{pmatrix}
\end{equation}

\textbf{Step 3 --- Compose the sample stack rotations.}
The total rotation applied to the sample is the ordered product of stage matrices,
from outermost (floor) to innermost (top):

\begin{equation}
  \mathbf{R}_\text{sample}
  = \mathbf{R}_\text{S2-1}(\theta_1)\,
    \mathbf{R}_\text{S2-2}(\theta_2)\,
    \mathbf{R}_\text{S2-3}(\theta_3)\,
    \mathbf{R}_\text{S2-4}(\theta_4)
\end{equation}

The order is significant: each stage rotates everything mounted above it.

\textbf{Step 4 --- Define \textbf{U}.}
The orientation matrix $\mathbf{U}$ maps a reference vector $\mathbf{h}$ expressed
in the crystal (sample) frame to its direction in the lab frame:

\begin{equation}
  \mathbf{h}_\text{lab} = \mathbf{R}_\text{sample}\,\mathbf{U}\,\mathbf{h}_\text{crystal}
\end{equation}

Given two non-parallel reference vectors measured at known stage angles, $\mathbf{U}$
is determined by:

\begin{equation}
  \mathbf{U}
  = \mathbf{R}_\text{sample}^{-1}\,
    \mathbf{M}_\text{lab}\,
    \mathbf{M}_\text{crystal}^{-1}
\end{equation}

where the columns of $\mathbf{M}_\text{lab}$ and $\mathbf{M}_\text{crystal}$ are the
reference vectors in the lab and crystal frames respectively, with the third column
taken as the cross product of the first two to complete the basis.

%-----------------------------------------------------------------------
\section*{Problem 2}

\subsection*{Are Different Basis Assignments Possible?}

Yes. The assignment of $\{\hat{x}, \hat{y}, \hat{z}\}$ to the three physical directions
was a choice, not a necessity. Any right-handed assignment is valid. For three physical
directions there are \textbf{6 possible right-handed assignments}:

\begin{center}
\begin{tabular}{clll}
\toprule
Assignment & $\hat{x}$ & $\hat{y}$ & $\hat{z}$ \\
\midrule
1 (chosen) & longitudinal & lateral      & vertical     \\
2          & lateral      & vertical     & longitudinal \\
3          & vertical     & longitudinal & lateral      \\
4          & lateral      & longitudinal & vertical     \\
5          & longitudinal & vertical     & lateral      \\
6          & vertical     & lateral      & longitudinal \\
\bottomrule
\end{tabular}
\end{center}

Assignments 1--3 are cyclic (even) permutations of (longitudinal, lateral, vertical);
assignments 4--6 are anti-cyclic (odd) permutations. Both sets are right-handed.

\subsection*{Stage Axis Vectors Under Each Assignment}

The physical rotation axes are unchanged; only their coordinate expression varies.

\begin{center}
\begin{tabular}{lllllllll}
\toprule
Stage & Physical Axis & A1 & A2 & A3 & A4 & A5 & A6 \\
\midrule
S1-1 & vertical     & $+\hat{z}$ & $+\hat{x}$ & $+\hat{y}$ & $+\hat{z}$ & $+\hat{y}$ & $+\hat{x}$ \\
S1-2 & lateral      & $+\hat{y}$ & $+\hat{z}$ & $+\hat{x}$ & $+\hat{x}$ & $+\hat{z}$ & $+\hat{y}$ \\
S2-1 & vertical     & $+\hat{z}$ & $+\hat{x}$ & $+\hat{y}$ & $+\hat{z}$ & $+\hat{y}$ & $+\hat{x}$ \\
S2-2 & lateral      & $+\hat{y}$ & $+\hat{z}$ & $+\hat{x}$ & $+\hat{x}$ & $+\hat{z}$ & $+\hat{y}$ \\
S2-3 & longitudinal & $+\hat{x}$ & $+\hat{y}$ & $+\hat{z}$ & $+\hat{y}$ & $+\hat{x}$ & $+\hat{z}$ \\
S2-4 & vertical     & $+\hat{z}$ & $+\hat{x}$ & $+\hat{y}$ & $+\hat{z}$ & $+\hat{y}$ & $+\hat{x}$ \\
\bottomrule
\end{tabular}
\end{center}

\subsection*{Effect on \textbf{U}}

The matrix $\mathbf{U}$ does not change physically --- it encodes the same geometric
relationship regardless of basis choice. What changes is its numerical representation.

Let $\mathbf{P}$ be the $3 \times 3$ orthogonal permutation matrix that transforms
coordinates from one basis assignment to another, satisfying $\mathbf{P}^{-1} = \mathbf{P}^\top$.
Then the orientation matrix under the new assignment is:

\begin{equation}
  \mathbf{U}' = \mathbf{P}\,\mathbf{U}\,\mathbf{P}^\top
\end{equation}

This is a similarity transformation, which preserves:

\begin{itemize}
  \item The eigenvalues of $\mathbf{U}$ (same rotation angles),
  \item $\det(\mathbf{U}') = \det(\mathbf{U}) = +1$ (proper rotation),
  \item The physical meaning of the orientation.
\end{itemize}

Only the numerical entries of $\mathbf{U}$ differ between assignments. The choice of
basis is purely a convention.

\end{document}
